%%%%%%%%%%%%%%%%%%%%%%%%%%%%%%%%%%%
% Cover Letter + point-by-point response — RSC Advances (RSC)
% Fresh submission after Digital Discovery rejection (DD-ART-07-2026-000486, 27-Aug-2026).
% Prior review disclosed; itemised response to the three Digital Discovery referees appended.
%%%%%%%%%%%%%%%%%%%%%%%%%%%%%%%%%%%

\documentclass[12pt,a4paper]{letter}
\usepackage[left=2.5cm, right=2.5cm, top=2.5cm, bottom=2.5cm]{geometry}
\usepackage{mathptmx}
\usepackage[T1]{fontenc}
\usepackage{microtype}
\usepackage{hyperref}
\usepackage[detect-all=true]{siunitx}
\usepackage{enumitem}

\signature{
  Jean-Pierre Tchapet Njafa\\
  (on behalf of all authors)
}

\address{
  Department of Physics, Faculty of Science\\
  University of Yaound\'{e}~1\\
  P.O. Box 812, Yaound\'{e}, Cameroon\\
  \texttt{jean-pierre.tchapet@facsciences-uy1.cm}
}

\date{September 9, 2026}

\begin{document}

\begin{letter}{
  Editorial Office, \textit{RSC Advances}\\
  Royal Society of Chemistry\\
  Thomas Graham House, Science Park\\
  Milton Road, Cambridge CB4 0WF, UK
}

\opening{Dear Editor,}

We submit our manuscript \textbf{``Experimental Inter-Report Scatter Bounds
Machine-Learning Accuracy for TADF Singlet--Triplet Gaps---Structure and
Semi-Empirical Features Reach the Same Limit''} for consideration as a Full Paper in
\textit{RSC Advances}.

\textbf{What the paper reports.}
The singlet--triplet gap $\Delta E_\text{ST}$ is the photophysical observable that
gates thermally activated delayed fluorescence (TADF). We measure where the
limit on predicting that observable from cheap information lies---and find it is
set by the experimental data, not by the model. On one controlled set of
\num{231} donor--acceptor emitters (\num{212} Bemis--Murcko scaffolds) we compare
three levels of description against experiment: the 2D molecular structure, a
semi-empirical excited-state calculation (GFN2-xTB/sTDA natural transition
orbitals), and a range-separated TD-DFT reference (wB97X-D4/def2-SVP,
Tamm--Dancoff, \num{231}/\num{231} converged, \num{609} core-hours). Structure
alone reaches mean absolute error \num{0.091}--\num{0.100}~eV; the semi-empirical
excited state gives \num{0.096}~eV; the TD-DFT gap gives \num{0.099}~eV after
in-fold calibration. They converge because the binding term is the label itself:
independent measurements of the same compound scatter by \num{0.072}~eV RMS, so
the per-molecule median the model is trained on is known only to \num{0.049}~eV---
about half the model error. (Computed references carry a separate functional
dependence of up to \num{0.58}~eV between B3LYP and CAM-B3LYP; that bounds what a
computed gap can mean and is not added into the experimental bound.) A further,
unexpected result separates accuracy from ranking: under a source-paper
(DOI-grouped) split, the semi-empirical features retain rank correlation the
fingerprints lose ($\Delta\rho = \num{0.19}$, \qty{95}{\percent} CI \numrange{0.04}{0.34},
excluding zero narrowly and on a single split), so they order candidates across
laboratories even where they do not improve the predicted value. This is a photochemistry benchmark with a
physical-organic conclusion---what structural information does and does not encode
about the exchange splitting---not solely a machine-learning exercise.

\textbf{Contributions.}
(i)~A controlled measurement of feature--target leakage for a derived-energy target.
Its severity is known to depend on the model class; we quantify that here. Regressing
$\Delta E_\text{ST}$ on features
containing the $S_1$ and $T_1$ energies is circular
($\Delta E_\text{ST}\equiv E_{S_1}-E_{T_1}$): under cross-validation a linear model
given those energies recovers the target exactly ($R^2 = 1.000$) against $R^2 = 0.36$
without them, while an axis-aligned random forest on the identical features reports
only $R^2 = 0.59$ and clears cross-validation unremarked. That is how a tree-based
screening study can carry this defect without the score looking wrong.
(ii)~Structure alone predicts the \emph{experimental} $\Delta E_\text{ST}$ as
accurately as the semi-empirical excited state---a scaffold-validated benchmark with
a label-permutation null ($p = 0.001$), the two representations equal to within
\qty{0.017}{\electronvolt} on paired folds. A new finding is the
dissociation between accuracy and ranking under laboratory transfer: the
semi-empirical features retain rank correlation under a source-paper split where
fingerprints do not ($\Delta\rho = \num{0.19}$, \qty{95}{\percent} CI excluding zero
narrowly, on a single split).
(iii)~A measured bound on the label itself---the median gap is known only to
\num{0.049}~eV---against which a \num{609}-core-hour TD-DFT reference and a frozen
pretrained molecular-transformer embedding both fail to beat the structure-only model.
(iv)~A triage filter that enriches low-gap candidates by
\num{1.2}--\num{1.4}$\times$ over random on held-out labelled data---useful for
prioritisation, not a device-level or emission-spectrum predictor.

\textbf{Why this is new.}
The advance is conceptual, not a better error. That $\Delta E_\text{ST}$ is reported
inconsistently has been shown for individual emitters (Shuaib \textit{et al.},
\textit{J.~Mater.~Chem.~C} 2026, \textbf{14}, 4353); we turn it into a quantitative
inter-laboratory variability budget over \num{231} molecules and \num{173} source
papers (single-report scatter \num{0.072}~eV RMS, median-label precision
\num{0.049}~eV) and show it is the term that caps every predictor we test,
semi-empirical or TD-DFT. Two further findings we have not seen reported: under a
source-paper split a model's accuracy advantage vanishes while the excited-state
descriptors keep the rank order the fingerprints lose ($\Delta\rho = \num{0.19}$,
\qty{95}{\percent} CI excluding zero narrowly, on a single split); and the
feature--target leakage that inflates computed-gap models is fully exposed by a linear
learner and substantially masked by a random forest. The multi-scale ML screening
pipelines and the pretrained or multimodal TADF models of 2025--2026 make neither
point, because none runs an experimental source-paper split or a controlled leakage
comparison. The work is the prediction-side, statistical complement to the
experimental measurement-methodology literature.

\textbf{Prior consideration.}
An earlier version of this manuscript was reviewed at \textit{Digital Discovery}
(DD-ART-07-2026-000486) and declined on 27~August~2026. The three referee reports
were constructive and we have used them to strengthen the paper. The decline cited
insufficient novelty against existing TADF screening work; the paragraph above states
what is new and what is not. An itemised response follows this letter. The main additions since that version are: the
\num{609}-core-hour range-separated TD-DFT reference run (Referee~1); a full data
and code deposit on Zenodo, previously incomplete (Referee~2); a source-paper
split, the removal of the library shortlist, and a trimmed
inverted-gap discussion (Referee~3).

\textbf{Open science.}
The experimental corpus, descriptor tables, trained models, SHAP analysis, TD-DFT
outputs, and all Python code are openly available on Zenodo
(DOI: 10.5281/zenodo.17436069) under a CC-BY-4.0 licence, with a SHA-256 manifest.

\textbf{APC waiver.}
The corresponding author is affiliated with the University of Yaound\'{e}~1,
Cameroon, a Research4Life country; we are eligible for a full APC waiver under the
RSC's Research4Life arrangement and will provide documentation on request.

\textbf{Suggested reviewers.}
None is a co-author or collaborator:
\begin{itemize}[nosep]
  \item Prof.\ Heather J.\ Kulik, Massachusetts Institute of Technology, USA
  (\texttt{hjkulik@mit.edu}) --- ML for molecular design
  \item Prof.\ Andrew P.\ Monkman, Durham University, UK
  (\texttt{a.p.monkman@durham.ac.uk}) --- TADF photophysics
  \item Prof.\ Olexandr Isayev, Carnegie Mellon University, USA
  (\texttt{olexandr@cmu.edu}) --- ML for chemistry
  \item Prof.\ Denis Jacquemin, Nantes Universit\'{e} / CNRS CEISAM, France
  (\texttt{Denis.Jacquemin@univ-nantes.fr}) --- TD-DFT excited-state benchmarks
\end{itemize}

The manuscript is not published elsewhere and is not under consideration by any
other journal. All authors have approved submission to \textit{RSC Advances}.

\closing{Yours sincerely,}

\end{letter}

\newpage

\begin{center}
  {\large\textbf{Response to the reviewers}}\\[2pt]
  Manuscript: ``Experimental Inter-Report Scatter Bounds Machine-Learning Accuracy
  for TADF Singlet--Triplet Gaps---Structure and Semi-Empirical Features Reach the
  Same Limit''\\
  Prior review: \textit{Digital Discovery} DD-ART-07-2026-000486
\end{center}

\medskip
\noindent
Referee comments are in italics. Every number cited below is reproducible from the
Zenodo deposit (DOI: 10.5281/zenodo.17436069). Section and table numbers refer to
the revised manuscript and its Supporting Information.

\bigskip
\noindent\textbf{\large Referee 1}

\medskip
\noindent\textbf{1.1}\quad\textit{``The dataset is tiny ($<$300 molecules), so it
is not surprising that it doesn't work.''}

\smallskip
\noindent\textbf{Response.}
Dataset size is the object of study here, not an oversight. A correlation-test
power analysis (Methods) shows that the observed Spearman $\rho = 0.36$ is
resolved at \num{80}\% power with $n \approx 60$; at $n = 231$ the achieved power
is $\approx 1$. We should be clear about the scope of that argument: it concerns the rank
correlation only, and does not speak to the width of the $R^2$ interval, which we address
separately and do not explain away.
More directly, we now report a threefold corpus expansion to \num{728}
experimental emitters through the same pipeline: the MAE \emph{rises} from
\num{0.098} to \num{0.124}~eV, $R^2$ falls from \num{0.18} to $-0.01$, and the
rank correlation roughly halves ($\rho$ \num{0.34} to \num{0.16})
(\S\,Ceilings; SI Table, expansion). Adding data makes the prediction worse, which
points at the label rather than the sample count---though we note in the manuscript that
the added molecules come through a different structure-resolution pipeline, so chemical
breadth and label quality are confounded in that particular test. The cleaner evidence is
direct: the model reaches \num{0.074}~eV on molecules whose replicate reports agree and
\num{0.106}~eV on those whose reports do not.

\medskip
\noindent\textbf{1.2}\quad\textit{``The correlation between experiment and theory
(DFT or semi-empirical) is rather good ($r > 0.95$). It is not difficult to
predict a computed singlet--triplet gap with ML and a large computed dataset.''}

\smallskip
\noindent\textbf{Response.}
We now test this claim directly. A range-separated TD-DFT reference gap
(wB97X-D4/def2-SVP, RIJCOSX, Tamm--Dancoff; ORCA~6; \texttt{nroots}~8, triplets
included; GFN2-xTB geometry) was computed for all \num{231} molecules
(\num{609} core-hours, \num{16.1}~min median per molecule; new Methods paragraph,
new SI subsection S4.3). Against \emph{experiment} the raw TD-DFT gap has Pearson
$r = 0.43$ and $R^2 = -16.3$: it is the worst predictor in the feature-comparison
table (MAE \num{0.654}~eV), because it carries a systematic $+\num{0.65}$~eV offset. The $r > 0.95$ figures in the
literature are theory-against-theory, or ML against a computed reference---the
circular set-up diagnosed in contribution~(i). Predicting the \emph{measured} gap
is the harder problem, and neither a semi-empirical nor a range-separated
excited-state calculation solves it here.

\medskip
\noindent\textbf{1.3}\quad\textit{``Particularly strange is the refusal to use the
computed gap as a descriptor while other computed quantities of similar cost are
used. The semi-empirical computed gap is probably much better than the ML gap
from structural parameters alone.''}

\smallskip
\noindent\textbf{Response.}
We now use the computed gap as a feature, at both levels of theory, and report
the result rather than omitting it (\S\,Benchmark; SI S4.3). The semi-empirical
(sTDA) gap as an added feature does not change the error (NTO model
\num{0.096} to \num{0.097}~eV; Morgan model unchanged). The TD-DFT gap as an
added feature lowers the Morgan model from \num{0.091} to \num{0.087}~eV---a
paired \num{95}\% CI of \numrange{0.0004}{0.0071}~eV and Wilcoxon $p = 0.016$, so
a real effect, but a \num{4}\% reduction bought with \num{609} core-hours. Used
alone, the semi-empirical gap is \emph{worse} than the structure-only forest
(calibrated MAE \num{0.104}~eV vs.\ \num{0.091}~eV), contrary to the referee's
expectation. We have also replaced the earlier assertion that the feature sets are
``statistically indistinguishable'' with a measured equivalence bound: paired on
identical folds, Morgan and NTO differ by \num{0.004}~eV with a \num{95}\% CI of
$[-0.017, +0.009]$~eV, so they are equivalent to within \num{0.017}~eV rather than
merely not separable. The energy scalars are excluded from the NTO descriptor set for a
specific reason---including $E_{S_1}$ and $E_{T_1}$ as features is the circularity
of contribution~(i); the spatial NTO overlap descriptors we do keep encode
exchange-integral physics without reducing to $E_{S_1}-E_{T_1}$.

\medskip
\noindent\textbf{1.4}\quad\textit{``I cannot see a practical use of this approach,
and the didactic value is limited.''}

\smallskip
\noindent\textbf{Response.}
The practical use is triage: on held-out labelled data the model enriches low-gap
emitters by \num{1.2}--\num{1.4}$\times$ over random selection---precision
\num{0.60} among the top \num{10} ranked candidates and \num{0.68} among the top
\num{50}, against a base rate of \num{0.49}. (These are counts of molecules, not
percentiles; the earlier submission conflated the two conventions and we have corrected
it.) We have re-scoped every claim to that regime and removed the projection onto the
unlabelled filtered library (see 3.4).

Two further results are directly usable. The first is a concrete recommendation: under a
source-paper split the semi-empirical NTO descriptors hold Spearman $\rho = \num{0.35}$
while fingerprints fall to \num{0.16}, so a group that must rank candidates against
measurements made in another laboratory should pay for the excited-state step, while a
group ranking within its own data need not. That is a decision the cost of those
descriptors can be weighed against.

The second is the didactic point, which we think is transferable. Under cross-validation a
linear model given the two constituent energies recovers the target exactly
($R^2 = \num{1.000}$) against $R^2 = \num{0.36}$ without them. More usefully, how severe
the leak \emph{looks} depends on the learner: a random forest is axis-aligned, does not
recover a difference of two continuous inputs exactly, and reaches only
$R^2 = \num{0.59}$ on identical leaked features. A tree-based study can therefore carry this defect while its
cross-validated score appears unremarkable, which may be why the failure mode persists.
Alongside it sits the calibrated chemistry ceiling: 2D structure and the semi-empirical
excited state land within \num{0.017}~eV of one another against experiment on paired
folds, and a calibrated TD-DFT gap within a further \num{0.008}~eV, while the raw
uncalibrated TD-DFT gap overestimates by \num{0.65}~eV.

\bigskip
\noindent\textbf{\large Referee 2}

\medskip
\noindent\textbf{2.1}\quad\textit{``Please provide a link to the Zenodo
repository, and make sure the actual data is there. The .zip I downloaded only
contained a README.md; the data was on GitHub. Please add it to Zenodo.''}

\smallskip
\noindent\textbf{Response.}
The Zenodo record 10.5281/zenodo.17436069 now holds the full deposit, not only
the README: the curated experimental corpus, the three descriptor tables, the
trained random-forest models, the per-molecule SHAP output, the new TD-DFT
reference outputs (\texttt{tddft\_reference\_231.csv} and timing), and every
analysis and figure-generation script, with a \texttt{manifest.csv} giving a
SHA-256 checksum for each file. The GitHub repository is retained as a mirror.
The DOI is cited in the abstract, Methods, Conclusions, and the data-availability
statement.

\bigskip
\noindent\textbf{\large Referee 3}

\medskip
\noindent\textbf{3.1}\quad\textit{``Predictive capacity is extremely weak and
statistically ambiguous: $R^2 = 0.25$--$0.31$ with \num{95}\% CIs spanning into
negative territory. A model with zero-crossing $R^2$ intervals cannot be
distinguished from a mean-value predictor. Spearman $\rho$ is insufficient to
claim predictive power.''}

\smallskip
\noindent\textbf{Response.}
We agree and have made this the framing rather than a caveat. The abstract now
states that the signal ``sits just above a mean-value baseline (MAE
\num{0.107}~eV)''; the Conclusions say the advantage is not resolved under the
harder split. What the signal \emph{is}: it clears a label-permutation null
($p = 0.001$); the Spearman $\rho = 0.36$ has a \num{95}\% CI of
\numrange{0.24}{0.47} that excludes zero; the paired MAE advantage over the
mean-value baseline excludes zero under scaffold CV. Actionable utility is
therefore triage only (\num{1.2}--\num{1.4}$\times$ enrichment on labelled data),
and we no longer make any molecule-level selection claim. Under the source-paper
(DOI-grouped) split introduced for this revision the advantage over the baseline
is \num{0.005}~eV with a CI that crosses zero---stated in the abstract, in
\S\,Benchmark, and as Limitations bullet~(6). One result cuts the other way and is now
reported: under that same split the semi-empirical NTO model holds Spearman
$\rho = \num{0.35}$ while the fingerprint model falls to \num{0.16}, a paired
difference of \num{0.19} with a \num{95}\% CI of $[0.04, 0.34]$ that excludes zero.
Ordering survives the change of protocol even where the predicted value does not, which
is what a triage filter is asked to do.

\medskip
\noindent\textbf{3.2}\quad\textit{``The dataset aggregates literature values
across unrecorded conditions (hosts, solvents, temperatures, techniques).
Compiling medians from uncontrolled environments degrades the target. Test or
discuss a solvent- and condition-controlled subset rather than blaming data
noise.''}

\smallskip
\noindent\textbf{Response.}
Methods now states explicitly that the corpus does not record measurement medium
per value, so the per-molecule median is a consensus over uncontrolled
conditions. We address the referee's test with the repeat-measured subset: for
the \num{102} molecules with independent reports, \num{58}\% have zero
inter-report spread, and this cleaner-label subset gives the study's
\emph{lowest} MAE, \num{0.080}~eV (\S\,Benchmark). We must be careful with that number
rather than offer it as support: a mean-value predictor restricted to the same
\num{102} molecules scores \num{0.089}~eV, so the model's advantage there is
\num{0.009}~eV with a confidence interval spanning zero, and the low absolute error
reflects the narrower spread of those labels rather than better prediction.
We tested the referee's proposition directly and it did not survive. An earlier
``clean-label'' subset of \num{208} molecules was mis-specified---it assigned zero spread
to the \num{129} compounds with no replicate at all, discarding only \num{23}.
Restricting properly to replicated molecules gives \num{0.074}~eV where reports agree and
\num{0.106}~eV where they do not, but a constant predictor fitted within each stratum
reproduces that ordering to within \num{0.001}~eV, because the strata differ in the
dispersion of the target. We therefore report no stratified evidence that label noise sets
the ceiling, and say so in the manuscript. On the referee's specific request we have gone further than stating a limitation. A
strictly solvent-controlled subset cannot be built: \texttt{phase} and
\texttt{atmosphere} are populated for \num{0} of the \num{1520} extracted rows. We
therefore ran the closest test the data permits, splitting on source paper, since values
drawn from a single publication share one laboratory's conditions. The model's advantage
over a constant predictor fitted within each group is \num{0.0071}~eV for the
\num{218} single-source molecules and \num{0.0024}~eV for the \num{13} multi-source
ones, a difference of \num{0.005}~eV whose \num{95}\% interval includes zero. The
result is null and we report it as such, noting that \num{13} molecules give the test
little power. We have also, throughout, compared every subset against a constant
predictor fitted within that subset, having found that an earlier subset argument of ours
was reproduced entirely by target dispersion.

\medskip
\noindent\textbf{3.3}\quad\textit{``With \num{212} scaffolds for \num{231}
molecules, \num{92.9}\% of scaffolds are singletons, so the GroupKFold scaffold
split is essentially a random split. On expansion to \num{728} molecules $R^2$
collapses to $-0.01$ and MAE worsens to \num{0.124}~eV, showing the model does
not generalise. The triage claim must be re-evaluated.''}

\smallskip
\noindent\textbf{Response.}
The singleton fraction is now stated in Methods and the SI, and we have added a
source-paper (DOI-grouped) \texttt{GroupKFold} over \num{173} papers as the
harder transfer test: it holds out a synthetic series and a laboratory's shared practice
together. Under it the Morgan model gives MAE \num{0.102}~eV ($R^2 = 0.18$) and
its advantage over the mean-value baseline is not statistically resolved
(\S\,Benchmark). The \num{728}-molecule expansion is promoted to its own Results
paragraph and read as the referee reads it---evidence against generalisation and
for a label ceiling. The triage claim is now explicitly restricted to held-out
labelled molecules, with the projection onto the filtered library removed.

\medskip
\noindent\textbf{3.4}\quad\textit{``The top-100 shortlist from the
\num{22194}-molecule library has a split-conformal \num{90}\% interval of
$\pm\num{0.18}$~eV, while the predicted gaps are \num{0.02}--\num{0.05}~eV. The
uncertainty engulfs the prediction; this offers no prioritisation value.''}

\smallskip
\noindent\textbf{Response.}
We agree and have removed the shortlist. No ranked candidate list is deposited.
The text (\S\,Library) now states that the split-conformal \num{90}\% interval
spans \num{0.36}~eV---about \num{7.6} times the predicted gap, a mean half-width of
\num{0.18}~eV---and that for every one of the \num{100} lowest-predicted molecules the
interval's lower bound is a negative (inverted) gap. Conformal prediction refuses
molecule-level selection at this accuracy; the filtered set is kept only to indicate
the scale of the chemical space the labelled-data triage would operate over. The
corresponding SI table has been deleted.

\medskip
\noindent\textbf{3.5}\quad\textit{``NTO descriptors are computed at ground-state
vertical geometries and miss excited-state relaxation and spin--orbit coupling.
Dismissing the \num{0.195}~eV adiabatic deviation as a rank-preserving offset is
inadequate; this likely explains why NTO features do not beat 2D fingerprints.''}

\smallskip
\noindent\textbf{Response.}
We have removed the dismissive wording. The \num{14}-molecule vertical-vs-adiabatic
comparison (CAM-B3LYP/def2-TZVP, MAE \num{0.195}~eV) is now reported as a bound on
the vertical approximation, not waved away, and the equivalence of NTO and 2D
features is now stated as a measured equivalence bound (\num{0.017}~eV on paired
folds) rather than as ``mechanistically informative''. Limitations bullet~(6) states
that the descriptors are ground-state vertical, carry no excited-state relaxation and
no spin--orbit coupling, and that this is a plausible part of why they do not
outperform 2D fingerprints. We agree that
missing excited-state relaxation is a plausible contributor to the NTO ceiling and
say so.

\medskip
\noindent\textbf{3.6}\quad\textit{``The single-determinant GFN2-xTB/sTDA framework
enforces Aufbau occupation and cannot capture the double-excitation character of
gap inversion; it predicts positive gaps for all inverted-gap candidates. The
INVEST narrative should be streamlined or removed.''}

\smallskip
\noindent\textbf{Response.}
Agreed, and the discussion is now short and factually corrected. A literature screen
flags five corpus compounds with an inverted report; for two of them (HAP-3MF,
TPA-NZP) the consensus median we train on is itself negative, and these are the only two
negative labels among the \num{231}. The model returns a positive gap for all five, so it
carries the wrong sign for both molecules whose own label is negative. We now state that
this is representational rather than statistical---the descriptors carry no signal that
separates an inverted-gap molecule from a conventional one, since the semi-empirical step
is computed under Aufbau ordering---and the ``expected falsification'' narrative is
removed. The earlier submission described these as
``external test cases'', which was wrong: they are inside the corpus.

\medskip
\noindent\textbf{3.7}\quad\textit{``An $R^2$ interval encompassing zero---does the
model provide actionable predictive utility for materials discovery?''}

\smallskip
\noindent\textbf{Response.}
For point prediction of $\Delta E_\text{ST}$: no, and the manuscript now says so
plainly. For triage: yes, within a stated bound---\num{1.2}--\num{1.4}$\times$
enrichment of low-gap emitters over random selection on held-out labelled data,
with conformal intervals that rule out single-molecule selection. The value of
the paper is the calibrated ceiling and the leakage diagnosis, both of which are
quantitative and reproducible.

\end{document}
